\documentclass[12pt,a4paper]{article}

% ============================================================
% Dan Koe "How to Fix Your Entire Life in 1 Day" 教學講義
% 通用版 - 適合一般民眾、學生、家庭閱讀
% 作者: 謝慕揚醫師 MD, PhD, FESC
% 來源: X (Twitter) @thedankoe
% 日期: 2026年1月13日
% ============================================================

% Drake_preamble.tex - LaTeX Preamble Template
% 作者: 謝慕揚, MD, PhD, FESC
% 
% 使用說明:
% 1. 表格請使用 longtable，不要有垂直分隔線 |
% 2. 表格內換行使用 \newline，不要使用 \\
% 3. 藥物名稱、檢驗、檢查請維持英文
% 4. Reference 格式請使用 \href 加上驗證過的 DOI

% Including essential packages for Chinese typesetting
\usepackage{xeCJK}
\usepackage{fontspec}
% Configuring Chinese font with Noto Serif CJK TC, fallback to SimSun
\IfFontExistsTF{Noto Serif CJK TC}{
  \setCJKmainfont{Noto Serif CJK TC}
  \setCJKsansfont{Noto Serif CJK TC}
  \setCJKmonofont{Noto Serif CJK TC}
}{\setCJKmainfont{SimSun}
  \setCJKsansfont{SimSun}
  \setCJKmonofont{SimSun}
}

% Including other necessary packages
\usepackage{geometry}
\usepackage{titlesec}
\usepackage{enumitem}
\usepackage{xcolor}
\usepackage{colortbl}
\usepackage{tcolorbox}
\usepackage{booktabs}
\usepackage{longtable}
\usepackage{array}
\usepackage{graphicx}
\usepackage{tikz}
\usepackage{fancyhdr}
\usepackage{multicol}
\usepackage{datetime2}
\usepackage{hyperref}
\usepackage{mdframed}
\usepackage{amsmath}
\usepackage{amssymb}
\usepackage{multirow}

\hypersetup{
    colorlinks=true,
    linkcolor=emergency_blue,
    citecolor=emergency_blue,
    urlcolor=emergency_blue,
    pdfborder={0 0 0}
}

% Defining page geometry
\geometry{
  top=2.5cm,
  bottom=2.5cm,
  left=2cm,
  right=2cm,
  headheight=25pt
}

% Adjusting topmargin to compensate for increased headheight
\addtolength{\topmargin}{-10pt}

% Defining colors
\definecolor{emergency_red}{RGB}{186,24,27}
\definecolor{emergency_blue}{RGB}{0,114,188}
\definecolor{emergency_gray}{RGB}{120,120,120}
\definecolor{highlight_yellow}{RGB}{255,250,205}

% Setting title formats
\titleformat{\section}[block]{\Large\bfseries\color{emergency_red}}{\thesection}{10pt}{}
\titleformat{\subsection}[block]{\large\bfseries\color{emergency_blue}}{\thesubsection}{8pt}{}
\titleformat{\subsubsection}[block]{\normalsize\bfseries\color{emergency_gray}}{\thesubsubsection}{6pt}{}

% Defining custom environment for important notes
\newmdenv[
    linecolor=emergency_red,
    backgroundcolor=highlight_yellow,
    linewidth=2pt,
    topline=true,
    bottomline=true,
    leftline=true,
    rightline=true,
    innertopmargin=10pt,
    innerbottommargin=10pt,
    innerleftmargin=10pt,
    innerrightmargin=10pt
]{importantbox}

% Defining note command with tcolorbox
\newcommand{\note}[1]{
    \par
    \noindent
    \begin{tcolorbox}[
        colback=emergency_blue!5!white,
        colframe=emergency_blue,
        title=\textbf{重點提醒},
        fonttitle=\bfseries,
        boxrule=0.5pt,
        arc=3pt,
        left=6pt,
        right=6pt,
        top=6pt,
        bottom=6pt
    ]
    #1
    \end{tcolorbox}
}

% Key findings box
\newcommand{\keyfinding}[1]{
    \par
    \noindent
    \begin{tcolorbox}[
        colback=yellow!10!white,
        colframe=orange!75!black,
        title=\textbf{核心發現摘要},
        fonttitle=\bfseries,
        boxrule=0.5pt,
        arc=3pt
    ]
    #1
    \end{tcolorbox}
}

% Defining custom date format for ISO 8601
\DTMnewdatestyle{iso}{
  \renewcommand{\DTMdisplaydate}[4]{\number##1-\number##2-\number##3}
  \renewcommand{\DTMDisplaydate}[4]{\number##1-\number##2-\number##3}
}
\DTMsetdatestyle{iso}
\newcommand{\mydate}{\DTMtoday}


% Custom header for this document
\pagestyle{fancy}
\fancyhf{}
\fancyhead[L]{\textcolor{emergency_red}{人生重啟指南}}
\fancyhead[R]{\textcolor{emergency_blue}{一天改變你的人生}}
\fancyfoot[C]{\textcolor{emergency_gray}{\thepage}}
\renewcommand{\headrulewidth}{1pt}
\renewcommand{\headrule}{\hbox to\headwidth{\color{emergency_red}\leaders\hrule height \headrulewidth\hfill}}

\begin{document}

% ============================================================
% 標題頁 - 通用版
% ============================================================

\begin{center}
{\LARGE\bfseries\textcolor{emergency_red}{人生重啟指南}}\\[0.3cm]
{\large 通用版：適合學生、家庭、所有想要改變的人}\\[0.5cm]
{\Large\textbf{How to Fix Your Entire Life in 1 Day}}\\[0.3cm]
{\large 如何在一天內重新定位你的人生}\\[0.5cm]
{\normalsize 原作者：Dan Koe (@thedankoe)}\\[0.2cm]
{\normalsize \href{https://x.com/thedankoe/status/2010751592346030461}{原文連結：X Post 2026年1月13日}}\\[0.3cm]
{\normalsize 整理：謝慕揚 MD, PhD, FESC \quad 日期：\mydate}
\end{center}

\vspace{0.5cm}

% ============================================================
% 核心訊息摘要
% ============================================================

\begin{tcolorbox}[colback=yellow!10!white,colframe=orange!75!black,title=\textbf{這篇文章在說什麼？}]
\textbf{三個改變人生的關鍵想法：}
\begin{enumerate}[leftmargin=*]
    \item \textbf{先改變「你是誰」，行為自然會跟著改變}——不是靠意志力硬撐
    \item \textbf{你做的每件事都有目的}——包括拖延、逃避，背後都有你沒察覺的原因
    \item \textbf{把人生當成遊戲來玩}——設定目標、每日任務、遊戲規則
\end{enumerate}

\textbf{適合誰閱讀？}
\begin{itemize}
    \item 想要改變但總是半途而廢的人
    \item 想幫助孩子建立好習慣的父母
    \item 想要一起成長的情侶或夫妻
    \item 正在迷茫的學生
\end{itemize}
\end{tcolorbox}

\tableofcontents
\newpage

% ============================================================
\section{為什麼新年目標總是失敗？}
% ============================================================

你有沒有這樣的經驗？

\begin{itemize}
    \item 新年說要減肥，堅持兩週就放棄
    \item 說要早起讀書，沒幾天就打回原形
    \item 說要存錢，但月底還是月光族
    \item 說要學英文，App 下載了就沒打開過
\end{itemize}

Dan Koe 說：\textbf{這不是因為你不夠自律，而是你用錯方法了。}

\subsection{兩種改變的方式}

\begin{longtable}{p{6.5cm}p{7.5cm}}
\toprule
\textbf{大多數人的做法（很難成功）} & \textbf{真正有效的做法} \\
\midrule
\endfirsthead

\bottomrule
\endfoot

「我要每天運動」\newline → 靠意志力硬撐 & 「我是一個重視健康的人」\newline → 運動變成自然的事 \\
\midrule
「我要少滑手機」\newline → 每次都忍不住 & 「我是一個珍惜時間的人」\newline → 自然想做更有意義的事 \\
\midrule
「我要認真讀書」\newline → 考完試就忘 & 「我是一個熱愛學習的人」\newline → 讀書變成享受 \\
\bottomrule
\end{longtable}

\begin{tcolorbox}[colback=emergency_blue!5!white, colframe=emergency_blue, title=\textbf{簡單來說}]
不要只改變你「做什麼」，要改變你認為「自己是誰」。

當你真的相信自己是某種人，你會自然而然做出符合那個身份的行為。
\end{tcolorbox}

% ============================================================
\section{第一個大觀念：成為對的人}
% ============================================================

\subsection{健身達人不是靠意志力}

想像一個身材很好的健身達人。你覺得他每天要「硬撐」才能吃得健康嗎？

答案是：不會。對他來說，吃垃圾食物才需要「硬撐」。

為什麼？因為他已經把「健康飲食」變成了他是誰的一部分。他不是在「減肥」，他就是一個「注重健康的人」。

\subsection{你現在的習慣，反映了你認為自己是誰}

\begin{longtable}{p{5cm}p{9cm}}
\toprule
\textbf{如果你常常...} & \textbf{可能是因為你認為自己是...} \\
\midrule
\endfirsthead

\bottomrule
\endfoot

拖延作業 & 「反正我不是讀書的料」 \\
\midrule
不敢舉手發言 & 「我不是那種會表現的人」 \\
\midrule
一直滑手機 & 「我就是需要放鬆，不然會累」 \\
\midrule
不敢嘗試新事物 & 「我不是那種敢冒險的人」 \\
\midrule
總是遲到 & 「我天生就是慢郎中」 \\
\bottomrule
\end{longtable}

\note{
\textbf{好消息是：}你可以選擇重新定義「你是誰」。這不是一夜之間的事，但從今天開始改變想法，行為就會慢慢跟上。
}

% ============================================================
\section{第二個大觀念：你做的每件事都有目的}
% ============================================================

這個觀念可能有點難接受，但很重要：

\begin{importantbox}
\textbf{你做的每一件事，包括你想要停止的壞習慣，背後都有一個你可能沒察覺的目的。}
\end{importantbox}

\subsection{拖延不是因為懶}

很多人以為拖延是因為懶惰或缺乏紀律。但 Dan Koe 說，拖延其實是在「達成某個目標」。

\textbf{例子：}
\begin{itemize}
    \item 你拖延寫報告 → 可能是因為「怕寫不好被批評」，所以乾脆不開始
    \item 你拖延整理房間 → 可能是因為「不想面對混亂的生活」
    \item 你拖延找工作 → 可能是因為「害怕被拒絕」
\end{itemize}

\subsection{常見的「隱藏目的」}

\begin{longtable}{p{4.5cm}p{4cm}p{5.5cm}}
\toprule
\textbf{表面行為} & \textbf{你以為的原因} & \textbf{可能的隱藏目的} \\
\midrule
\endfirsthead

\multicolumn{3}{c}{\textbf{表格續接}} \\
\toprule
\textbf{表面行為} & \textbf{你以為的原因} & \textbf{可能的隱藏目的} \\
\midrule
\endhead

\bottomrule
\endfoot

考試前不讀書 & 我太懶了 & 如果沒認真讀還考不好，\newline 就不是我「笨」 \\
\midrule
不敢表白 & 我害羞 & 不被拒絕就不會受傷 \\
\midrule
不嘗試新事物 & 我沒興趣 & 不嘗試就不會失敗 \\
\midrule
總是說「隨便」 & 我沒意見 & 不表態就不用負責 \\
\midrule
不運動 & 我沒時間 & 維持現狀比較舒適 \\
\bottomrule
\end{longtable}

\begin{tcolorbox}[colback=emergency_blue!5!white, colframe=emergency_blue, title=\textbf{這代表什麼？}]
要改變行為，你需要先察覺背後的「隱藏目的」，然後問自己：

「這個目的真的值得嗎？它帶給我的安全感，比我想要的人生更重要嗎？」
\end{tcolorbox}

% ============================================================
\section{第三個大觀念：恐懼如何塑造你}
% ============================================================

\subsection{你的很多想法不是你自己想出來的}

從小到大，我們吸收了很多來自父母、老師、朋友、社會的想法：

\begin{itemize}
    \item 「穩定的工作最重要」
    \item 「不要太出風頭」
    \item 「先求有，再求好」
    \item 「你又不是那塊料」
\end{itemize}

這些想法可能是對的，也可能只是別人的恐懼傳給你的。

\subsection{身份認同是怎麼形成的？}

\begin{enumerate}
    \item 你想要達成一個目標（例如：被大人稱讚）
    \item 你發現某些行為可以達成這個目標
    \item 你重複這些行為，變成習慣
    \item 這些習慣變成「你是誰」的一部分
    \item 你開始保護這個身份，即使它不再適合你
\end{enumerate}

\note{
\textbf{例子：}小時候你發現「乖乖聽話」可以得到稱讚，所以你變成一個「聽話的孩子」。長大後，你可能還在習慣性地聽從別人，即使那不是你真正想要的。
}

\subsection{當身份被挑戰時會發生什麼？}

當有人質疑你相信的事情，你會感到不舒服，甚至想反擊。這是因為：

\textbf{攻擊你的想法 = 攻擊你是誰}

這就是為什麼改變這麼難——改變意味著「承認過去的自己是錯的」，這會讓人不舒服。

% ============================================================
\section{第四個大觀念：心智發展的階段}
% ============================================================

人的心智會隨著年齡和經歷成長。了解這些階段可以幫助你知道自己在哪裡，以及可以往哪裡走。

\subsection{簡化版心智發展階段}

\begin{longtable}{clp{8cm}}
\toprule
\textbf{階段} & \textbf{名稱} & \textbf{特徵（用簡單的話說）} \\
\midrule
\endfirsthead

\multicolumn{3}{c}{\textbf{表格續接}} \\
\toprule
\textbf{階段} & \textbf{名稱} & \textbf{特徵} \\
\midrule
\endhead

\bottomrule
\endfoot

1 & 衝動 & 想到什麼就做什麼，沒有考慮後果 \\
\midrule
2 & 自我保護 & 學會保護自己，知道什麼話該說、什麼話不該說 \\
\midrule
3 & 從眾 & 跟著群體走，大家怎麼做我就怎麼做 \\
\midrule
4 & 自我覺察 & 開始發現「我好像跟別人想的不一樣」 \\
\midrule
5 & 有原則 & 建立自己的價值觀，知道什麼對自己重要 \\
\midrule
6 & 獨立思考 & 了解自己的想法也是被環境塑造的，開始更開放 \\
\midrule
7 & 系統思考 & 能看到事情的全貌，知道自己也有盲點 \\
\midrule
8 & 框架覺察 & 知道所有的想法和身份都是「人造的」，可以選擇 \\
\midrule
9 & 統一 & 不再執著於「成為什麼」，享受當下 \\
\bottomrule
\end{longtable}

\begin{tcolorbox}[colback=emergency_blue!5!white, colframe=emergency_blue, title=\textbf{你可能在哪裡？}]
大多數青少年在第3-4階段（從眾→自我覺察）。\\
大多數成年人在第4-5階段（自我覺察→有原則）。\\
能到第6階段以上的人比較少，但任何人都可以透過學習和反思進步。
\end{tcolorbox}

% ============================================================
\section{第五個大觀念：什麼是真正的聰明？}
% ============================================================

Dan Koe 引用了一個很酷的定義：

\begin{importantbox}
\textbf{「聰明的唯一真正測試是：你有沒有得到你想要的人生。」}\\
\hfill — Naval Ravikant
\end{importantbox}

這不是說考試成績不重要，而是說：真正的聰明是「能夠達成自己的目標」。

\subsection{聰明人做什麼不一樣？}

\begin{longtable}{p{7cm}p{7cm}}
\toprule
\textbf{比較不聰明的做法} & \textbf{比較聰明的做法} \\
\midrule
\endfirsthead

\bottomrule
\endfoot

遇到困難就放棄 & 遇到困難就想辦法繞過去 \\
\midrule
一直犯同樣的錯 & 從錯誤中學習，下次做不同的事 \\
\midrule
只會抱怨 & 想辦法改變能改變的 \\
\midrule
覺得「我就是這樣」 & 相信自己可以成長 \\
\midrule
害怕嘗試新事物 & 把失敗當成學習的機會 \\
\bottomrule
\end{longtable}

\subsection{成功的公式}

\begin{center}
\fbox{\Large \textbf{成功 = 主動 + 機會 + 學習能力}}
\end{center}

\begin{itemize}
    \item \textbf{主動}：願意採取行動，不是光想不做
    \item \textbf{機會}：現在的網路時代，機會比以前多很多
    \item \textbf{學習能力}：能從經驗中學習，不斷調整方向
\end{itemize}

% ============================================================
\section{第六個大觀念：一天重啟你的人生}
% ============================================================

Dan Koe 設計了一個「一日重啟協議」，幫助你在一天內重新思考人生方向。

\subsection{改變的三個階段}

\begin{enumerate}
    \item \textbf{不滿（Dissonance）}：對現狀感到厭煩，覺得「不能再這樣下去了」
    \item \textbf{迷茫（Uncertainty）}：不知道接下來該怎麼辦
    \item \textbf{發現（Discovery）}：找到新方向，開始快速進步
\end{enumerate}

這個協議的目標是幫助你走過這三個階段。

\subsection{早上：想清楚你「不要」什麼}

花15-30分鐘思考這些問題：

\begin{enumerate}
    \item 你生活中有什麼讓你不開心，但你一直忍耐著？
    \item 你最常抱怨什麼？（但從來沒有真的去改變）
    \item 如果未來五年什麼都不變，你的生活會是什麼樣子？
    \item 想像十年後的自己——如果你從來沒有改變，你會後悔什麼？
\end{enumerate}

\subsection{白天：打斷自動駕駛}

在手機上設定提醒，全天問自己這些問題：

\begin{itemize}
    \item \textbf{11:00}：我現在做的事，是在幫助我還是在逃避什麼？
    \item \textbf{14:00}：如果有人看我過去兩小時的行為，他們會覺得我想要什麼？
    \item \textbf{17:00}：今天我有沒有做什麼讓我離目標更近？
    \item \textbf{20:00}：今天什麼時候我最有活力？什麼時候最沒勁？
\end{itemize}

\subsection{晚上：設定新方向}

\begin{enumerate}
    \item 寫下一句話：「我拒絕讓我的人生變成\_\_\_\_\_\_\_\_\_\_」
    \item 寫下一句話：「我想要成為\_\_\_\_\_\_\_\_\_\_的人」
    \item 一年後，什麼必須改變，我才會知道自己在進步？
    \item 明天，我可以做什麼小事來開始這個改變？
\end{enumerate}

% ============================================================
\section{第七個大觀念：把人生當遊戲玩}
% ============================================================

遊戲為什麼讓人上癮？因為它們有：
\begin{itemize}
    \item 清楚的目標（打倒魔王）
    \item 明確的規則（不能作弊）
    \item 即時的回饋（升級、獲得裝備）
    \item 適當的挑戰（不會太難也不會太簡單）
\end{itemize}

\textbf{如果我們把人生也設計成這樣呢？}

\subsection{人生遊戲設定}

\begin{longtable}{p{4cm}p{4cm}p{6cm}}
\toprule
\textbf{遊戲元素} & \textbf{人生對應} & \textbf{你的設定} \\
\midrule
\endfirsthead

\bottomrule
\endfoot

\textbf{獲勝條件} & 願景 & 你想要的理想人生是什麼？ \\
\midrule
\textbf{失敗結果} & 反願景 & 你絕對不想要的人生是什麼？ \\
\midrule
\textbf{主線任務} & 一年目標 & 一年後你想達成什麼？ \\
\midrule
\textbf{Boss戰} & 月計畫 & 這個月要完成什麼挑戰？ \\
\midrule
\textbf{每日任務} & 日常行動 & 每天做什麼可以推進目標？ \\
\midrule
\textbf{遊戲規則} & 限制條件 & 有什麼是你不願意犧牲的？ \\
\bottomrule
\end{longtable}

\begin{tcolorbox}[colback=emergency_blue!5!white, colframe=emergency_blue, title=\textbf{為什麼這樣做有用？}]
當你有清楚的目標和每日任務，你的大腦會更容易專注。

就像玩遊戲時你不會想看手機，因為你知道自己在做什麼、要往哪裡去。
\end{tcolorbox}

% ============================================================
\section{【實作練習】我的人生重啟工作表}
% ============================================================

\textit{找一個安靜的時間，拿出紙筆，認真填寫這些問題。}

% ------------------------------------------------------------
\subsection{Part 1：認識現在的自己}
% ------------------------------------------------------------

\textbf{1. 我對自己最不滿意的三件事：}
\begin{center}
\fbox{\parbox{0.9\textwidth}{
\begin{enumerate}
    \item \vspace{1.2cm}
    \item \vspace{1.2cm}
    \item \vspace{1.2cm}
\end{enumerate}
}}
\end{center}

\textbf{2. 我最常用什麼藉口？}
\begin{center}
\fbox{\parbox{0.9\textwidth}{
例：「我沒時間」「我做不到」「等以後再說」\\[0.5cm]
我的藉口：\vspace{2cm}
}}
\end{center}

\textbf{3. 如果有人觀察我的行為（不是聽我說的話），他們會認為我想要什麼？}
\begin{center}
\fbox{\parbox{0.9\textwidth}{\vspace{2.5cm}}}
\end{center}

\textbf{4. 我現在的身份宣言是什麼？（誠實地寫）}
\begin{center}
\fbox{\parbox{0.9\textwidth}{
「我是那種 \underline{\hspace{8cm}} 的人」
\vspace{1cm}
}}
\end{center}

\newpage

% ------------------------------------------------------------
\subsection{Part 2：想像未來}
% ------------------------------------------------------------

\textbf{5. 如果什麼都不改變，五年後的我會是什麼樣子？}
\begin{center}
\fbox{\parbox{0.9\textwidth}{\vspace{4cm}}}
\end{center}

\textbf{6. 我真正想要的生活是什麼樣子？（不要管「現實」，先大膽想像）}
\begin{center}
\fbox{\parbox{0.9\textwidth}{\vspace{4cm}}}
\end{center}

\textbf{7. 我想要成為什麼樣的人？}
\begin{center}
\fbox{\parbox{0.9\textwidth}{
\LARGE 「我想成為那種 \underline{\hspace{5cm}} 的人」
\vspace{1cm}
}}
\end{center}

\newpage

% ------------------------------------------------------------
\subsection{Part 3：制定行動計畫}
% ------------------------------------------------------------

\textbf{我的人生遊戲設定：}

\begin{longtable}{p{4cm}p{10cm}}
\toprule
\textbf{設定項目} & \textbf{我的答案} \\
\midrule
\endfirsthead

\bottomrule
\endfoot

\textbf{反願景}\newline {\small 我絕對不要的人生} & \fbox{\parbox{8.5cm}{\vspace{2.5cm}}} \\
\midrule
\textbf{願景}\newline {\small 我想要的人生} & \fbox{\parbox{8.5cm}{\vspace{2.5cm}}} \\
\midrule
\textbf{一年目標}\newline {\small 一年後我要達成...} & \fbox{\parbox{8.5cm}{\vspace{2cm}}} \\
\midrule
\textbf{這個月的挑戰}\newline {\small 我要完成...} & \fbox{\parbox{8.5cm}{\vspace{2cm}}} \\
\midrule
\textbf{每天的小行動}\newline {\small 我每天會做...} & \fbox{\parbox{8.5cm}{\vspace{2cm}}} \\
\midrule
\textbf{我的底線}\newline {\small 我不願意犧牲...} & \fbox{\parbox{8.5cm}{\vspace{2cm}}} \\
\bottomrule
\end{longtable}

% ============================================================
\section{給不同讀者的建議}
% ============================================================

% ------------------------------------------------------------
\subsection{給學生的話}
% ------------------------------------------------------------

\textbf{你可能覺得：}「我還是學生，這些跟我有什麼關係？」

\textbf{其實：}學生時期是建立身份認同最重要的時候。你現在相信的「我是什麼樣的人」，會影響你未來幾十年。

\begin{tcolorbox}[colback=yellow!10!white,colframe=orange!75!black,title=\textbf{學生可以嘗試的事}]
\begin{enumerate}
    \item \textbf{換一個身份宣言}
    \begin{itemize}
        \item 從「我數學不好」→「我正在學習數學」
        \item 從「我很害羞」→「我正在練習表達自己」
    \end{itemize}

    \item \textbf{問自己：「我逃避這件事的真正原因是什麼？」}
    \begin{itemize}
        \item 不是「我懶」，而是「我怕什麼？」
    \end{itemize}

    \item \textbf{設定一個「每日任務」}
    \begin{itemize}
        \item 例如：每天讀10頁書、每天運動15分鐘
        \item 重點是「每天」，不是「量多」
    \end{itemize}
\end{enumerate}
\end{tcolorbox}

% ------------------------------------------------------------
\subsection{給父母和孩子一起讀}
% ------------------------------------------------------------

這篇文章的觀念也可以幫助親子關係：

\textbf{給父母：}
\begin{itemize}
    \item 孩子的很多「問題行為」背後都有一個目的——可能是想要關注、安全感、或逃避壓力
    \item 與其說「你怎麼這麼懶」，不如問「你在擔心什麼？」
    \item 幫助孩子建立正面的身份認同：「你是一個認真的孩子」比「你要更認真」更有效
\end{itemize}

\textbf{給孩子：}
\begin{itemize}
    \item 你現在相信的「我是什麼樣的人」不是永遠固定的
    \item 你可以選擇成為你想成為的人
    \item 失敗不代表你笨，只代表你正在學習
\end{itemize}

\begin{tcolorbox}[colback=emergency_blue!5!white, colframe=emergency_blue, title=\textbf{親子對話題目}]
\begin{enumerate}
    \item 「你覺得你是什麼樣的人？」「你想成為什麼樣的人？」
    \item 「有什麼事情你想做但一直沒做？你在擔心什麼？」
    \item 「如果你可以改變一件事，你會改變什麼？」
    \item 「我們可以一起設定一個目標嗎？」
\end{enumerate}
\end{tcolorbox}

% ------------------------------------------------------------
\subsection{給情侶和夫妻}
% ------------------------------------------------------------

這篇文章的觀念也可以應用在關係中：

\textbf{一起成長：}
\begin{itemize}
    \item 分享彼此的「願景」和「反願景」
    \item 討論：「我們想要什麼樣的生活？」
    \item 設定「我們的」目標，而不只是「我的」目標
\end{itemize}

\textbf{理解彼此：}
\begin{itemize}
    \item 對方的「問題行為」背後可能有你沒看到的恐懼或需求
    \item 與其批評，不如好奇：「你為什麼這樣做？你在擔心什麼？」
\end{itemize}

\begin{tcolorbox}[colback=yellow!10!white,colframe=orange!75!black,title=\textbf{情侶/夫妻對話題目}]
\begin{enumerate}
    \item 「你理想中我們五年後的生活是什麼樣子？」
    \item 「有什麼事情你一直想做但沒告訴我？」
    \item 「我們可以一起設定一個目標嗎？」
    \item 「有什麼是我可以做的，讓你更容易達成你的目標？」
\end{enumerate}
\end{tcolorbox}

% ============================================================
\section{結語：改變從今天開始}
% ============================================================

\begin{enumerate}
    \item \textcolor{red}{\textbf{改變從身份開始}}：不是「我要做什麼」，而是「我要成為誰」
    \item \textcolor{red}{\textbf{察覺隱藏目的}}：你的每個行為背後都有原因，找到它
    \item \textcolor{red}{\textbf{設定你的遊戲}}：有目標、有任務、有規則的人生更容易專注
    \item \textcolor{red}{\textbf{行動比完美重要}}：不需要等到準備好，今天就開始小小的改變
\end{enumerate}

\vspace{0.5cm}

\begin{center}
\fbox{\parbox{0.8\textwidth}{
\centering
\textbf{你的第一步}\\[0.3cm]
今天就花15分鐘填寫「實作練習」工作表。\\
不需要完美，只需要開始。\\[0.3cm]
改變不是一夜之間發生的，\\
但每一個小步驟都在讓你離想要的人生更近。
}}
\end{center}

% ============================================================
\section{參考資料}
% ============================================================

\begin{enumerate}
    \item Koe D. How to fix your entire life in 1 day. \href{https://x.com/thedankoe/status/2010751592346030461}{X (Twitter). 2026年1月13日.}

    \item Csikszentmihalyi M. \textit{Flow: The Psychology of Optimal Experience}. 1990.（心流理論）

    \item Dweck C. \textit{Mindset: The New Psychology of Success}. 2006.（成長心態）

    \item Clear J. \textit{Atomic Habits}. 2018.（原子習慣）
\end{enumerate}

% ============================================================
\section{名詞解釋}
% ============================================================

\begin{longtable}{p{4cm}p{10cm}}
\toprule
\textbf{名詞} & \textbf{簡單解釋} \\
\midrule
\endfirsthead

\bottomrule
\endfoot

身份認同 & 你認為「自己是什麼樣的人」 \\
\midrule
願景 & 你想要的理想未來 \\
\midrule
反願景 & 你絕對不想要的未來 \\
\midrule
隱藏目的 & 你做某件事的真正原因（可能你自己都沒發現） \\
\midrule
自動駕駛 & 不用思考就照習慣做事 \\
\midrule
心流 & 完全專注在一件事上，忘記時間的感覺 \\
\bottomrule
\end{longtable}

\end{document}
